\section*{Resumen}
\addcontentsline{toc}{section}{Resumen}

Se hizo \textit{fine-tuning} de BERT-base y DistilBERT-base sobre tres tareas
de clasificación de texto (SST-2, AG News y Yelp Polarity) bajo un protocolo
idéntico, y se comparó a los dos modelos en desempeño y en eficiencia. Sobre
esa base se ejecutó un estudio de ablación de seis configuraciones que varía
la cabeza de clasificación y qué parte del \textit{transformer} se entrena.

Los resultados sostienen tres afirmaciones. Primera: con el \SI{61}{\percent}
de los parámetros, DistilBERT conserva entre el \SI{98,9}{\percent} y el
\SI{99,5}{\percent} de la \textit{accuracy} de BERT, y la brecha depende de la
tarea --- se abre en análisis de sentimiento sobre frases cortas y se cierra
hasta desaparecer en clasificación de temas. Segunda: medida en condiciones
comparables, la inferencia de DistilBERT es exactamente \best{2$\times$} más
rápida en los nueve puntos de una rejilla de \textit{batch} y longitud, que es
lo que predice pasar de 12 capas a 6. Tercera: la arquitectura de la cabeza de
clasificación es irrelevante --- cinco configuraciones distintas caben en
\num{0,0014} de F1 --- mientras que congelar el \textit{encoder} cuesta entre
\num{3,4} y \num{8,4} puntos de F1 según la tarea.

En total son \num{24} corridas de entrenamiento, unas \num{2,5} horas de GPU
en una NVIDIA RTX A6000. Todas las métricas de este informe se generan
automáticamente desde los archivos JSON versionados en el repositorio.
